% =============================================================================
% def.tex — Centralised notation and macro definitions
% Include via % =============================================================================
% def.tex — Centralised notation and macro definitions
% Include via % =============================================================================
% def.tex — Centralised notation and macro definitions
% Include via % =============================================================================
% def.tex — Centralised notation and macro definitions
% Include via \input{def} in your main .tex file.
% Change a symbol once here and it updates everywhere.
% =============================================================================

% ---- Packages required by the macros below ----
\usepackage{amsmath,amssymb,amsthm}
\usepackage{bm}              % bold math
\usepackage{mathtools}        % \coloneqq, \DeclarePairedDelimiter, etc.
\usepackage{algorithm}
\usepackage{algpseudocode}
\usepackage{enumitem}

% ---- Theorem-like environments (requires amsthm loaded above) ----
\newtheorem{theorem}{Theorem}
\newtheorem{lemma}[theorem]{Lemma}
\newtheorem{proposition}[theorem]{Proposition}
\newtheorem{corollary}[theorem]{Corollary}
\newtheorem{definition}{Definition}
\newtheorem{remark}{Remark}
\newtheorem{assumption}{Assumption}

\renewcommand{\algorithmicrequire}{\textbf{Input:}}
\renewcommand{\algorithmicensure}{\textbf{Output:}}


\newcommand{\cM}{\mathcal{M}}
\newcommand{\cN}{\mathcal{N}}
\newcommand{\saddle}{{\vc{x}^{*}}}
\DeclareMathOperator{\grad}{grad}
\DeclareMathOperator{\Hess}{Hess}
\newcommand{\proj}[2]{\mathcal{P}_{#1}\!\left({#2}\right)}
\newcommand{\ind}{\textbf{Ind }}
\newcommand{\defs}{\coloneqq}
\newcommand{\saddleset}{\mathcal{S}}         % set of saddle points
\newcommand{\eigmean}{\bar{\lambda}}          % mean eigenvalue
\newcommand{\crit}[1]{e\left(#1\right)}
\newcommand{\tangent}[2]{T_{#1}#2}
\newcommand{\indicator}[1]{\mathbf{1}_{#1}}

% =============================================================================
% 1. Common math operators
% =============================================================================
\DeclareMathOperator*{\argmax}{arg\,max}
\DeclareMathOperator*{\argmin}{arg\,min}
\DeclareMathOperator{\softmax}{softmax}
\DeclareMathOperator{\relu}{ReLU}
\DeclareMathOperator{\diag}{diag}
\DeclareMathOperator{\tr}{tr}           % trace
\DeclareMathOperator{\rank}{rank}
\DeclareMathOperator{\sign}{sign}

% =============================================================================
% 2. Scalar / vector / matrix / tensor conventions
%    Change the styling once and every occurrence follows.
% =============================================================================
\newcommand{\sca}[1]{#1}                         % scalar  (italic, default)
\newcommand{\vc}[1]{\bm{#1}}                     % vector  (bold italic)
\newcommand{\mat}[1]{\bm{#1}}                    % matrix  (bold italic)
\newcommand{\ten}[1]{\bm{\mathcal{#1}}}          % tensor  (bold calligraphic)

% =============================================================================
% 3. Named vectors / matrices you use often
%    e.g. \x, \W — rename freely to match your paper's notation.
% =============================================================================
% -- Vectors --
\newcommand{\x}{{\vc{x}}}          % input
\newcommand{\y}{\vc{y}}          % target / output
\newcommand{\z}{\vc{z}}          % latent
\newcommand{\h}{\vc{h}}          % hidden state
\newcommand{\e}{\vc{e}}          % embedding / error

% -- Matrices --
\newcommand{\W}{\mat{W}}         % weight matrix
\newcommand{\V}{\mat{V}}         % weight matrix
\newcommand{\A}{\mat{A}}         % generic matrix A
\newcommand{\I}{\mat{I}}         % identity matrix

% -- Special scalars --
\newcommand{\loss}{\mathcal{L}}  % loss function
\newcommand{\obj}{\mathcal{J}}   % objective function
\newcommand{\reg}{\Omega}        % regulariser

% =============================================================================
% 4. Sets and spaces
% =============================================================================
\newcommand{\RR}{\mathbb{R}}      % real numbers
\newcommand{\NN}{\mathbb{N}}      % natural numbers
\newcommand{\ZZ}{\mathbb{Z}}      % integers

\newcommand{\CC}{\mathbb{C}}      % integers
\newcommand{\XX}{\mathcal{X}}     % input space
\newcommand{\YY}{\mathcal{Y}}     % output space
\newcommand{\DD}{\mathcal{D}}     % dataset
\newcommand{\spr}{\mathbb{S}}     % sphere

% =============================================================================
% 5. Probability / statistics
% =============================================================================
\newcommand{\E}{\mathbb{E}}                      % expectation
\newcommand{\Var}{\mathrm{Var}}                   % variance
\newcommand{\Cov}{\mathrm{Cov}}                   % covariance
\newcommand{\KL}{\mathrm{KL}}                     % KL divergence
\newcommand{\prob}[1]{p\!\left(#1\right)}         % p(.)
\newcommand{\given}{\mid}                          % conditional bar

% =============================================================================
% 6. Delimiters (auto-sized)
% =============================================================================
\DeclarePairedDelimiter{\abs}{\lvert}{\rvert}      % |x|
\DeclarePairedDelimiter{\norm}{\lVert}{\rVert}     % ||x||
\DeclarePairedDelimiter{\inner}{\langle}{\rangle}  % <x, y>
\DeclarePairedDelimiter{\floor}{\lfloor}{\rfloor}
\DeclarePairedDelimiter{\ceil}{\lceil}{\rceil}

% =============================================================================
% 7. Convenience shortcuts
% =============================================================================
% \newcommand{\T}{^{\mkern-1.5mu\mathsf{T}}}        % transpose
% \newcommand{\inv}{^{-1}}                            % inverse
% % \newcommand{\dd}{\mathrm{d}}                        % differential d
% \newcommand{\pd}[2]{\frac{\partial #1}{\partial #2}} % partial derivative
% \newcommand{\grad}{\nabla}                           % gradient symbol

% % =============================================================================
% % 8. Colours for the paper (easy to rebrand)
% % =============================================================================
% \definecolor{OurBlue}{HTML}{1F77B4}
% \definecolor{OurOrange}{HTML}{FF7F0E}
% \definecolor{OurGreen}{HTML}{2CA02C}
% \definecolor{OurRed}{HTML}{D62728}

% % =============================================================================
% % 9. Model / method names (change once, updates everywhere)
% % =============================================================================
\newcommand{\ours}{\textsc{YourMethod}}            % your method name
 in your main .tex file.
% Change a symbol once here and it updates everywhere.
% =============================================================================

% ---- Packages required by the macros below ----
\usepackage{amsmath,amssymb,amsthm}
\usepackage{bm}              % bold math
\usepackage{mathtools}        % \coloneqq, \DeclarePairedDelimiter, etc.
\usepackage{algorithm}
\usepackage{algpseudocode}
\usepackage{enumitem}

% ---- Theorem-like environments (requires amsthm loaded above) ----
\newtheorem{theorem}{Theorem}
\newtheorem{lemma}[theorem]{Lemma}
\newtheorem{proposition}[theorem]{Proposition}
\newtheorem{corollary}[theorem]{Corollary}
\newtheorem{definition}{Definition}
\newtheorem{remark}{Remark}
\newtheorem{assumption}{Assumption}

\renewcommand{\algorithmicrequire}{\textbf{Input:}}
\renewcommand{\algorithmicensure}{\textbf{Output:}}


\newcommand{\cM}{\mathcal{M}}
\newcommand{\cN}{\mathcal{N}}
\newcommand{\saddle}{{\vc{x}^{*}}}
\DeclareMathOperator{\grad}{grad}
\DeclareMathOperator{\Hess}{Hess}
\newcommand{\proj}[2]{\mathcal{P}_{#1}\!\left({#2}\right)}
\newcommand{\ind}{\textbf{Ind }}
\newcommand{\defs}{\coloneqq}
\newcommand{\saddleset}{\mathcal{S}}         % set of saddle points
\newcommand{\eigmean}{\bar{\lambda}}          % mean eigenvalue
\newcommand{\crit}[1]{e\left(#1\right)}
\newcommand{\tangent}[2]{T_{#1}#2}
\newcommand{\indicator}[1]{\mathbf{1}_{#1}}

% =============================================================================
% 1. Common math operators
% =============================================================================
\DeclareMathOperator*{\argmax}{arg\,max}
\DeclareMathOperator*{\argmin}{arg\,min}
\DeclareMathOperator{\softmax}{softmax}
\DeclareMathOperator{\relu}{ReLU}
\DeclareMathOperator{\diag}{diag}
\DeclareMathOperator{\tr}{tr}           % trace
\DeclareMathOperator{\rank}{rank}
\DeclareMathOperator{\sign}{sign}

% =============================================================================
% 2. Scalar / vector / matrix / tensor conventions
%    Change the styling once and every occurrence follows.
% =============================================================================
\newcommand{\sca}[1]{#1}                         % scalar  (italic, default)
\newcommand{\vc}[1]{\bm{#1}}                     % vector  (bold italic)
\newcommand{\mat}[1]{\bm{#1}}                    % matrix  (bold italic)
\newcommand{\ten}[1]{\bm{\mathcal{#1}}}          % tensor  (bold calligraphic)

% =============================================================================
% 3. Named vectors / matrices you use often
%    e.g. \x, \W — rename freely to match your paper's notation.
% =============================================================================
% -- Vectors --
\newcommand{\x}{{\vc{x}}}          % input
\newcommand{\y}{\vc{y}}          % target / output
\newcommand{\z}{\vc{z}}          % latent
\newcommand{\h}{\vc{h}}          % hidden state
\newcommand{\e}{\vc{e}}          % embedding / error

% -- Matrices --
\newcommand{\W}{\mat{W}}         % weight matrix
\newcommand{\V}{\mat{V}}         % weight matrix
\newcommand{\A}{\mat{A}}         % generic matrix A
\newcommand{\I}{\mat{I}}         % identity matrix

% -- Special scalars --
\newcommand{\loss}{\mathcal{L}}  % loss function
\newcommand{\obj}{\mathcal{J}}   % objective function
\newcommand{\reg}{\Omega}        % regulariser

% =============================================================================
% 4. Sets and spaces
% =============================================================================
\newcommand{\RR}{\mathbb{R}}      % real numbers
\newcommand{\NN}{\mathbb{N}}      % natural numbers
\newcommand{\ZZ}{\mathbb{Z}}      % integers

\newcommand{\CC}{\mathbb{C}}      % integers
\newcommand{\XX}{\mathcal{X}}     % input space
\newcommand{\YY}{\mathcal{Y}}     % output space
\newcommand{\DD}{\mathcal{D}}     % dataset
\newcommand{\spr}{\mathbb{S}}     % sphere

% =============================================================================
% 5. Probability / statistics
% =============================================================================
\newcommand{\E}{\mathbb{E}}                      % expectation
\newcommand{\Var}{\mathrm{Var}}                   % variance
\newcommand{\Cov}{\mathrm{Cov}}                   % covariance
\newcommand{\KL}{\mathrm{KL}}                     % KL divergence
\newcommand{\prob}[1]{p\!\left(#1\right)}         % p(.)
\newcommand{\given}{\mid}                          % conditional bar

% =============================================================================
% 6. Delimiters (auto-sized)
% =============================================================================
\DeclarePairedDelimiter{\abs}{\lvert}{\rvert}      % |x|
\DeclarePairedDelimiter{\norm}{\lVert}{\rVert}     % ||x||
\DeclarePairedDelimiter{\inner}{\langle}{\rangle}  % <x, y>
\DeclarePairedDelimiter{\floor}{\lfloor}{\rfloor}
\DeclarePairedDelimiter{\ceil}{\lceil}{\rceil}

% =============================================================================
% 7. Convenience shortcuts
% =============================================================================
% \newcommand{\T}{^{\mkern-1.5mu\mathsf{T}}}        % transpose
% \newcommand{\inv}{^{-1}}                            % inverse
% % \newcommand{\dd}{\mathrm{d}}                        % differential d
% \newcommand{\pd}[2]{\frac{\partial #1}{\partial #2}} % partial derivative
% \newcommand{\grad}{\nabla}                           % gradient symbol

% % =============================================================================
% % 8. Colours for the paper (easy to rebrand)
% % =============================================================================
% \definecolor{OurBlue}{HTML}{1F77B4}
% \definecolor{OurOrange}{HTML}{FF7F0E}
% \definecolor{OurGreen}{HTML}{2CA02C}
% \definecolor{OurRed}{HTML}{D62728}

% % =============================================================================
% % 9. Model / method names (change once, updates everywhere)
% % =============================================================================
\newcommand{\ours}{\textsc{YourMethod}}            % your method name
 in your main .tex file.
% Change a symbol once here and it updates everywhere.
% =============================================================================

% ---- Packages required by the macros below ----
\usepackage{amsmath,amssymb,amsthm}
\usepackage{bm}              % bold math
\usepackage{mathtools}        % \coloneqq, \DeclarePairedDelimiter, etc.
\usepackage{algorithm}
\usepackage{algpseudocode}
\usepackage{enumitem}

% ---- Theorem-like environments (requires amsthm loaded above) ----
\newtheorem{theorem}{Theorem}
\newtheorem{lemma}[theorem]{Lemma}
\newtheorem{proposition}[theorem]{Proposition}
\newtheorem{corollary}[theorem]{Corollary}
\newtheorem{definition}{Definition}
\newtheorem{remark}{Remark}
\newtheorem{assumption}{Assumption}

\renewcommand{\algorithmicrequire}{\textbf{Input:}}
\renewcommand{\algorithmicensure}{\textbf{Output:}}


\newcommand{\cM}{\mathcal{M}}
\newcommand{\cN}{\mathcal{N}}
\newcommand{\saddle}{{\vc{x}^{*}}}
\DeclareMathOperator{\grad}{grad}
\DeclareMathOperator{\Hess}{Hess}
\newcommand{\proj}[2]{\mathcal{P}_{#1}\!\left({#2}\right)}
\newcommand{\ind}{\textbf{Ind }}
\newcommand{\defs}{\coloneqq}
\newcommand{\saddleset}{\mathcal{S}}         % set of saddle points
\newcommand{\eigmean}{\bar{\lambda}}          % mean eigenvalue
\newcommand{\crit}[1]{e\left(#1\right)}
\newcommand{\tangent}[2]{T_{#1}#2}
\newcommand{\indicator}[1]{\mathbf{1}_{#1}}

% =============================================================================
% 1. Common math operators
% =============================================================================
\DeclareMathOperator*{\argmax}{arg\,max}
\DeclareMathOperator*{\argmin}{arg\,min}
\DeclareMathOperator{\softmax}{softmax}
\DeclareMathOperator{\relu}{ReLU}
\DeclareMathOperator{\diag}{diag}
\DeclareMathOperator{\tr}{tr}           % trace
\DeclareMathOperator{\rank}{rank}
\DeclareMathOperator{\sign}{sign}

% =============================================================================
% 2. Scalar / vector / matrix / tensor conventions
%    Change the styling once and every occurrence follows.
% =============================================================================
\newcommand{\sca}[1]{#1}                         % scalar  (italic, default)
\newcommand{\vc}[1]{\bm{#1}}                     % vector  (bold italic)
\newcommand{\mat}[1]{\bm{#1}}                    % matrix  (bold italic)
\newcommand{\ten}[1]{\bm{\mathcal{#1}}}          % tensor  (bold calligraphic)

% =============================================================================
% 3. Named vectors / matrices you use often
%    e.g. \x, \W — rename freely to match your paper's notation.
% =============================================================================
% -- Vectors --
\newcommand{\x}{{\vc{x}}}          % input
\newcommand{\y}{\vc{y}}          % target / output
\newcommand{\z}{\vc{z}}          % latent
\newcommand{\h}{\vc{h}}          % hidden state
\newcommand{\e}{\vc{e}}          % embedding / error

% -- Matrices --
\newcommand{\W}{\mat{W}}         % weight matrix
\newcommand{\V}{\mat{V}}         % weight matrix
\newcommand{\A}{\mat{A}}         % generic matrix A
\newcommand{\I}{\mat{I}}         % identity matrix

% -- Special scalars --
\newcommand{\loss}{\mathcal{L}}  % loss function
\newcommand{\obj}{\mathcal{J}}   % objective function
\newcommand{\reg}{\Omega}        % regulariser

% =============================================================================
% 4. Sets and spaces
% =============================================================================
\newcommand{\RR}{\mathbb{R}}      % real numbers
\newcommand{\NN}{\mathbb{N}}      % natural numbers
\newcommand{\ZZ}{\mathbb{Z}}      % integers

\newcommand{\CC}{\mathbb{C}}      % integers
\newcommand{\XX}{\mathcal{X}}     % input space
\newcommand{\YY}{\mathcal{Y}}     % output space
\newcommand{\DD}{\mathcal{D}}     % dataset
\newcommand{\spr}{\mathbb{S}}     % sphere

% =============================================================================
% 5. Probability / statistics
% =============================================================================
\newcommand{\E}{\mathbb{E}}                      % expectation
\newcommand{\Var}{\mathrm{Var}}                   % variance
\newcommand{\Cov}{\mathrm{Cov}}                   % covariance
\newcommand{\KL}{\mathrm{KL}}                     % KL divergence
\newcommand{\prob}[1]{p\!\left(#1\right)}         % p(.)
\newcommand{\given}{\mid}                          % conditional bar

% =============================================================================
% 6. Delimiters (auto-sized)
% =============================================================================
\DeclarePairedDelimiter{\abs}{\lvert}{\rvert}      % |x|
\DeclarePairedDelimiter{\norm}{\lVert}{\rVert}     % ||x||
\DeclarePairedDelimiter{\inner}{\langle}{\rangle}  % <x, y>
\DeclarePairedDelimiter{\floor}{\lfloor}{\rfloor}
\DeclarePairedDelimiter{\ceil}{\lceil}{\rceil}

% =============================================================================
% 7. Convenience shortcuts
% =============================================================================
% \newcommand{\T}{^{\mkern-1.5mu\mathsf{T}}}        % transpose
% \newcommand{\inv}{^{-1}}                            % inverse
% % \newcommand{\dd}{\mathrm{d}}                        % differential d
% \newcommand{\pd}[2]{\frac{\partial #1}{\partial #2}} % partial derivative
% \newcommand{\grad}{\nabla}                           % gradient symbol

% % =============================================================================
% % 8. Colours for the paper (easy to rebrand)
% % =============================================================================
% \definecolor{OurBlue}{HTML}{1F77B4}
% \definecolor{OurOrange}{HTML}{FF7F0E}
% \definecolor{OurGreen}{HTML}{2CA02C}
% \definecolor{OurRed}{HTML}{D62728}

% % =============================================================================
% % 9. Model / method names (change once, updates everywhere)
% % =============================================================================
\newcommand{\ours}{\textsc{YourMethod}}            % your method name
 in your main .tex file.
% Change a symbol once here and it updates everywhere.
% =============================================================================

% ---- Packages required by the macros below ----
\usepackage{amsmath,amssymb,amsthm}
\usepackage{bm}              % bold math
\usepackage{mathtools}        % \coloneqq, \DeclarePairedDelimiter, etc.
\usepackage{algorithm}
\usepackage{algpseudocode}
\usepackage{enumitem}

% ---- Theorem-like environments (requires amsthm loaded above) ----
\newtheorem{theorem}{Theorem}
\newtheorem{lemma}[theorem]{Lemma}
\newtheorem{proposition}[theorem]{Proposition}
\newtheorem{corollary}[theorem]{Corollary}
\newtheorem{definition}{Definition}
\newtheorem{remark}{Remark}
\newtheorem{assumption}{Assumption}

\renewcommand{\algorithmicrequire}{\textbf{Input:}}
\renewcommand{\algorithmicensure}{\textbf{Output:}}


\newcommand{\cM}{\mathcal{M}}
\newcommand{\cN}{\mathcal{N}}
\newcommand{\saddle}{{\vc{x}^{*}}}
\DeclareMathOperator{\grad}{grad}
\DeclareMathOperator{\Hess}{Hess}
\newcommand{\proj}[2]{\mathcal{P}_{#1}\!\left({#2}\right)}
\newcommand{\ind}{\textbf{Ind }}
\newcommand{\defs}{\coloneqq}
\newcommand{\saddleset}{\mathcal{S}}         % set of saddle points
\newcommand{\eigmean}{\bar{\lambda}}          % mean eigenvalue
\newcommand{\crit}[1]{e\left(#1\right)}
\newcommand{\tangent}[2]{T_{#1}#2}
\newcommand{\indicator}[1]{\mathbf{1}_{#1}}

% =============================================================================
% 1. Common math operators
% =============================================================================
\DeclareMathOperator*{\argmax}{arg\,max}
\DeclareMathOperator*{\argmin}{arg\,min}
\DeclareMathOperator{\softmax}{softmax}
\DeclareMathOperator{\relu}{ReLU}
\DeclareMathOperator{\diag}{diag}
\DeclareMathOperator{\tr}{tr}           % trace
\DeclareMathOperator{\rank}{rank}
\DeclareMathOperator{\sign}{sign}

% =============================================================================
% 2. Scalar / vector / matrix / tensor conventions
%    Change the styling once and every occurrence follows.
% =============================================================================
\newcommand{\sca}[1]{#1}                         % scalar  (italic, default)
\newcommand{\vc}[1]{\bm{#1}}                     % vector  (bold italic)
\newcommand{\mat}[1]{\bm{#1}}                    % matrix  (bold italic)
\newcommand{\ten}[1]{\bm{\mathcal{#1}}}          % tensor  (bold calligraphic)

% =============================================================================
% 3. Named vectors / matrices you use often
%    e.g. \x, \W — rename freely to match your paper's notation.
% =============================================================================
% -- Vectors --
\newcommand{\x}{{\vc{x}}}          % input
\newcommand{\y}{\vc{y}}          % target / output
\newcommand{\z}{\vc{z}}          % latent
\newcommand{\h}{\vc{h}}          % hidden state
\newcommand{\e}{\vc{e}}          % embedding / error

% -- Matrices --
\newcommand{\W}{\mat{W}}         % weight matrix
\newcommand{\V}{\mat{V}}         % weight matrix
\newcommand{\A}{\mat{A}}         % generic matrix A
\newcommand{\I}{\mat{I}}         % identity matrix

% -- Special scalars --
\newcommand{\loss}{\mathcal{L}}  % loss function
\newcommand{\obj}{\mathcal{J}}   % objective function
\newcommand{\reg}{\Omega}        % regulariser

% =============================================================================
% 4. Sets and spaces
% =============================================================================
\newcommand{\RR}{\mathbb{R}}      % real numbers
\newcommand{\NN}{\mathbb{N}}      % natural numbers
\newcommand{\ZZ}{\mathbb{Z}}      % integers

\newcommand{\CC}{\mathbb{C}}      % integers
\newcommand{\XX}{\mathcal{X}}     % input space
\newcommand{\YY}{\mathcal{Y}}     % output space
\newcommand{\DD}{\mathcal{D}}     % dataset
\newcommand{\spr}{\mathbb{S}}     % sphere

% =============================================================================
% 5. Probability / statistics
% =============================================================================
\newcommand{\E}{\mathbb{E}}                      % expectation
\newcommand{\Var}{\mathrm{Var}}                   % variance
\newcommand{\Cov}{\mathrm{Cov}}                   % covariance
\newcommand{\KL}{\mathrm{KL}}                     % KL divergence
\newcommand{\prob}[1]{p\!\left(#1\right)}         % p(.)
\newcommand{\given}{\mid}                          % conditional bar

% =============================================================================
% 6. Delimiters (auto-sized)
% =============================================================================
\DeclarePairedDelimiter{\abs}{\lvert}{\rvert}      % |x|
\DeclarePairedDelimiter{\norm}{\lVert}{\rVert}     % ||x||
\DeclarePairedDelimiter{\inner}{\langle}{\rangle}  % <x, y>
\DeclarePairedDelimiter{\floor}{\lfloor}{\rfloor}
\DeclarePairedDelimiter{\ceil}{\lceil}{\rceil}

% =============================================================================
% 7. Convenience shortcuts
% =============================================================================
% \newcommand{\T}{^{\mkern-1.5mu\mathsf{T}}}        % transpose
% \newcommand{\inv}{^{-1}}                            % inverse
% % \newcommand{\dd}{\mathrm{d}}                        % differential d
% \newcommand{\pd}[2]{\frac{\partial #1}{\partial #2}} % partial derivative
% \newcommand{\grad}{\nabla}                           % gradient symbol

% % =============================================================================
% % 8. Colours for the paper (easy to rebrand)
% % =============================================================================
% \definecolor{OurBlue}{HTML}{1F77B4}
% \definecolor{OurOrange}{HTML}{FF7F0E}
% \definecolor{OurGreen}{HTML}{2CA02C}
% \definecolor{OurRed}{HTML}{D62728}

% % =============================================================================
% % 9. Model / method names (change once, updates everywhere)
% % =============================================================================
\newcommand{\ours}{\textsc{YourMethod}}            % your method name
